\section{Konsekwencje dla modeli: niestandardowe liczby i paradoks Skolema}\label{sec:modele-konsekwencje}
\subsection*{Te same zdania nie muszą opisywać jednego świata}\label{sec:15:1}
Izomorfizm dwóch struktur to bijekcja ich dziedzin zachowująca stałe, działania i relacje. Struktury izomorficzne różnią się tylko nazwami elementów. Modele o różnych mocach nie mogą być izomorficzne, bo nie istnieje nawet bijekcja ich dziedzin.

Wniosek~\ref{res:14:10} ma więc bardzo konkretny sens: jeżeli aksjomaty pierwszego rzędu w przeliczalnym języku dopuszczają jakiś nieskończony świat, to nie opisują jednoznacznie jego wielkości. Te same aksjomaty dopuszczają światy przeliczalne i nieprzeliczalne.

Dotyczy to także teorii \emph{zupełnej}, która rozstrzyga każde zdanie swojego języka. Zupełność odpowiada na pytanie „czy ustalono odpowiedź na każde pytanie pierwszego rzędu?”, a jedyność modelu z dokładnością do izomorfizmu --- „czy ustalono całą strukturę?”. To różne wymagania. Zupełna teoria może mieć nieizomorficzne modele zgadzające się co do wszystkich zdań, które ta teoria potrafi wypowiedzieć.

Nie chodzi tylko o brak wystarczająco sprytnego \emph{skończonego} zestawu aksjomatów. Nawet dodanie wszystkich zdań pierwszego rzędu prawdziwych w danej nieskończonej strukturze nie usunie ograniczenia dotyczącego mocy modeli. Twierdzenie nie wymaga bowiem efektywności ani skończoności zbioru aksjomatów.

\subsection*{Niestandardowa arytmetyka: liczba większa od każdego numerału}\label{sec:15:2}
Weźmy $\operatorname{Th}(\mathbb N)$, czyli \emph{wszystkie} zdania arytmetyczne prawdziwe w standardowych liczbach naturalnych. To najsilniejszy możliwy opis $\mathbb N$ za pomocą prawdziwych zdań tego języka. Mimo to nie wyznacza on $\mathbb N$ jednoznacznie, nawet wśród modeli przeliczalnych. Ten ostatni fakt wymaga zwartości, nie wynika z samego istnienia modeli różnych mocy.

\begin{stwierdzenie}[Przeliczalny niestandardowy model prawdziwej arytmetyki]\label{res:15:1}
Istnieje przeliczalny model $\mathcal M\models\operatorname{Th}(\mathbb N)$ zawierający element $c^{\mathcal M}$ większy od interpretacji każdego standardowego numerału $\overline n$. Model ten nie jest izomorficzny ze standardową strukturą $\mathbb N$.
\end{stwierdzenie}
\begin{proof}
Dodajemy nową stałą $c$ i rozpatrujemy zbiór zdań
\[
\Gamma=\operatorname{Th}(\mathbb N)\cup\{\overline n<c:n\in\mathbb N\}.
\]
Każdy skończony fragment ma model: bierzemy zwykłe $\mathbb N$, a $c$ interpretujemy jako liczbę większą od wszystkich skończenie wielu numerałów występujących w dodatkowych nierównościach. Z Lematu~\ref{res:14:7} istnieje model całego $\Gamma$. Twierdzenie~\ref{res:14:5}, zastosowane do języka z $c$, daje jego przeliczalny podmodel elementarny, zachowujący także element nazwany $c$.

Po pominięciu symbolu $c$ nadal mamy strukturę arytmetyczną oraz wyróżniony element jej dziedziny. Każdy izomorfizm arytmetyczny zachowuje $0$ i $S$, więc przenosi interpretację $\overline n$ na zwykłą liczbę $n$. Obraz $c^{\mathcal M}$ musiałby być liczbą naturalną większą od każdego $n\in\mathbb N$, a takiej liczby w standardowym $\mathbb N$ nie ma.
\end{proof}

\textbf{Jak to sobie wyobrażać?} Model zawiera elementy $0,S(0),S(S(0)),\ldots$, ale jego dziedzina nie kończy się na tej zewnętrznie opisanej liście. Element $c^{\mathcal M}$ nie jest żadnym ze standardowych numerałów. Nie znaleźliśmy nowej zwykłej liczby naturalnej „za wszystkimi innymi”; skonstruowaliśmy \emph{inną strukturę}, w której symbole arytmetyczne mają interpretacje spełniające te same zdania pierwszego rzędu.

\textbf{Dlaczego indukcja tego nie zabrania?} W PA indukcja jest schematem dla własności wyrażalnych formułą, także z parametrami. Nie mówi o każdym zewnętrznie wybranym podzbiorze dziedziny modelu. Zbiór elementów będących interpretacjami standardowych numerałów zawiera $0$ i jest zamknięty ze względu na następnik, ale w niestandardowym modelu nie jest definiowalny formułą arytmetyczną, nawet z parametrami. Gdyby był, odpowiednia instancja indukcji wymusiłaby, że jest całą dziedziną, wbrew istnieniu $c^{\mathcal M}$. Dlatego schemat indukcji i istnienie niestandardowych elementów nie przeczą sobie.

\subsection*{Paradoks Skolema: przeliczalny model wierzący w nieprzeliczalność}\label{sec:15:3}
Załóżmy, że istnieje model teorii mnogości ZFC. ZFC jest sformułowana w przeliczalnym języku pierwszego rzędu, którego jedynym symbolem pozalogicznym jest relacja przynależności $\in$. Jej modele są nieskończone, więc Twierdzenie~\ref{res:14:5} daje model przeliczalny. Jest to wniosek \emph{warunkowy}: nie dowodzi istnienia modelu ZFC ani niesprzeczności ZFC bez żadnych założeń.

ZFC dowodzi jednak, że zbiór liczb rzeczywistych jest nieprzeliczalny. Czy dostaliśmy sprzeczność: „w przeliczalnym świecie jest nieprzeliczalny zbiór”? Nie, bo porównujemy dwa różne znaczenia słów „istnieje funkcja”.

Niech $M$ będzie przeliczalną dziedziną takiego modelu, a $r\in M$ obiektem, który model uważa za zbiór liczb rzeczywistych. Zewnętrznie możemy rozpatrzyć zbiór
\[
E_r=\{a\in M:a\mathrel{\in^{\mathcal M}}r\}.
\]
Jest on co najwyżej przeliczalny, ponieważ jest podzbiorem $M$. Zapis $\in^{\mathcal M}$ oznacza \emph{relację interpretowaną w modelu}; nie zakładamy, że jest ona zwykłą przynależnością między naszymi zbiorami.

Wewnętrzne zdanie „liczby rzeczywiste są nieprzeliczalne” mówi natomiast, że \emph{nie ma w modelu obiektu}, który model uznaje za bijekcję między swoim zbiorem liczb naturalnych a $r$. Zewnętrzne wyliczenie $E_r$ nie musi być reprezentowane przez obiekt modelu. Co więcej, jego dziedziną jest nasze standardowe $\mathbb N$, podczas gdy wewnętrzna funkcja miałaby dziedzinę będącą zbiorem liczb naturalnych w sensie modelu. Samo zewnętrzne wyliczenie nie dostarcza więc zakazanej funkcji \emph{wewnątrz modelu}.

\begin{quote}
Model zawiera wystarczająco dużo obiektów, aby spełniać swoje aksjomaty, ale nie musi zawierać wszystkich funkcji i podzbiorów, które my potrafimy rozpatrzyć, patrząc na jego dziedzinę z zewnątrz.
\end{quote}

Dotyczy to również aksjomatu zbioru potęgowego: jego kwantyfikatory przebiegają po obiektach modelu, nie po wszystkich zewnętrznych podzbiorach odpowiedniego zbioru $E_r$. To wyjaśnienie nazywa się rozwiązaniem \emph{paradoksu Skolema}. Nie obala ono twierdzenia Cantora i nie pokazuje, że nasz rzeczywisty zbiór $\mathbb R$ jest przeliczalny.

\subsection*{A gdyby użyć logiki drugiego rzędu?}\label{sec:15:4}
Pełna semantyka drugiego rzędu pozwala kwantyfikować po \emph{wszystkich} podzbiorach dziedziny. Drugorzędowy aksjomat indukcji obejmuje wtedy także zewnętrzny zbiór
\[
A=\{(S^{\mathcal M})^n(0^{\mathcal M}):n\in\mathbb N\}.
\]
Zbiór ten zawiera $0^{\mathcal M}$ i jest zamknięty na $S^{\mathcal M}$, więc pełna indukcja drugiego rzędu daje $A=M$. Odwzorowanie $n\mapsto(S^{\mathcal M})^n(0^{\mathcal M})$ jest zatem surjekcją. Jest też różnowartościowe: równość wyrazów o indeksach $n<m$ po $n$ skróceniach następnika dawałaby $0^{\mathcal M}=(S^{\mathcal M})^{m-n}(0^{\mathcal M})$, wbrew temu, że zero nie jest następnikiem. Równania rekurencyjne dla $+$ i $\cdot$ pokazują, przez zewnętrzną indukcję po drugim argumencie, że ta bijekcja zachowuje działania. Definicja $<$ zapewnia zachowanie porządku. Otrzymujemy więc izomorfizm ze standardową strukturą $\mathbb N$, a nie tylko bijekcję dziedzin.

To nie jest kontrprzykład do twierdzenia Löwenheima--Skolema: zmieniliśmy logikę i znaczenie kwantyfikatorów. Zyskujemy silniejszy opis semantyczny, ale tracimy możliwość efektywnego, poprawnego i zupełnego rachunku dowodowego dla pełnej semantyki. W szczególności gdyby taki rachunek obejmował wszystkie konsekwencje pełnych aksjomatów drugorzędowej arytmetyki, wyliczałby wszystkie prawdy pierwszego rzędu o $\mathbb N$, wbrew Wnioskowi~\ref{res:9:5}. Przy semantyce Henkina, w której zmienne drugiego rzędu przebiegają tylko po wybranej rodzinie podzbiorów, ten argument o jednoznaczności już nie działa: rodzina może nie zawierać zewnętrznego zbioru numerałów.

